\documentclass[11pt,parskip=half-]{scrartcl}
\usepackage[headsepline,footsepline]{scrlayer-scrpage}
\usepackage{amsmath, amssymb}
\usepackage{fontspec}
\usepackage{newpxtext,newpxmath}
\usepackage{graphicx,grffile}
\usepackage{setspace}
\usepackage{booktabs}
\usepackage{hyperref}
\usepackage{float}
\usepackage{enumitem}
\usepackage{lastpage}
\usepackage[backend=biber,style=authoryear]{biblatex}

\addbibresource{clarity2_sap.bib} %Import the bibliography file

\renewcommand\pagemark{{\usekomafont{pagenumber}\thepage\ of \pageref{LastPage}}}

\hypersetup{
    colorlinks=true,
    linkcolor=blue,
    filecolor=magenta,      
    urlcolor=cyan,
    pdftitle={CLARITY 2 -  Statistical Analysis Plan},
    pdfpagemode=FullScreen,
    }

% Setup header
\automark*{section}
\clearpairofpagestyles
\pagestyle{scrheadings}
\ihead{\headmark}
\ifoot{CLARITY 2.0 SAP - Version 0.2}
\ofoot{\pagemark}

% Setup line spacing
\renewcommand{\arraystretch}{1}
\onehalfspacing
\AfterTOCHead{\singlespacing}

%\RedeclareSectionCommand[
%  afterskip=1sp]{section}
%\RedeclareSectionCommand[
%  afterskip=1sp]{subsection}
%\RedeclareSectionCommand[
%  afterskip=1sp]{subsubsection}

% Custom commands
\providecommand{\tightlist}{%
  \setlength{\itemsep}{0pt}\setlength{\parskip}{0pt}}

  % enumitem options
\setlist[itemize]{parsep=2pt}
\setlist[enumerate]{parsep=2pt}


\begin{document}
\title{Statistical Analysis Plan}
\subtitle{CLARITY 2.0}
\author{James Totterdell}
\date{\today}
\makeatletter
\begin{titlepage}
    \begin{center}
        \includegraphics[width=0.7\linewidth]{clarity-logo.jpg}\\[4ex]
        {\huge \bfseries  \@title }\\[2ex]
        {\LARGE \bfseries  \@subtitle }\\[2ex]
        {\large An Investigator Initiated, International Multi-Centre, Multi-Arm, Multi-Stage
        Randomised Double Blind Placebo Controlled Trial of Angiotensin Receptor
        Blocker (ARB) and Chemokine Receptor Type 2 (CCR2) Antagonist for the Treatment
        of COVID-19}\\[2ex]
        {\LARGE Version 0.2}\\[2ex]
        {\LARGE \@date}\\[10ex]
        \fbox{\begin{tabular}{p{0.20\linewidth}p{0.5\linewidth}}
                Sponsor:         & The University of Sydney \\
                Registration:    & CTC0312                  \\
                Study Co-Chairs: & Professor Meg Jardine    \\
                                 & Professor Vivekanand Jha
            \end{tabular}}
    \end{center}
\end{titlepage}
\makeatother
\thispagestyle{empty}
\newpage

\tableofcontents

\clearpage

\section*{Preface}

This statistical analysis plan (SAP) outlines the data and procedures for analysing effectiveness and safety of trial interventions from the protocol CLARITY 2.0: An Investigator Initiated, International Multi-Centre, Multi-Arm, Multi-Stage Randomised Double Blind Placebo Controlled Trial of Angiotensin Receptor Blocker (ARB) \& Chemokine Receptor Type 2 (CCR2) Antagonist for the Treatment of COVID-19.

The content of this SAP aims to satisfy the recommended items as outlined in \cite{gamble2017}.

The following documents were reviewed when preparing this SAP (most recent version only):

\begin{itemize}
    \item CLARITY 2.0 Study Protocol version 2.0 24 September 2021
\end{itemize}


The planned analyses are similar to those undertaken in the CLARITY trial \cite{hockham2021protocol, mcgree2021controlled}.

\section*{Version History}

\subsection*{Note}

On 2022-07-25 we were informed that the funder had issued a letter of termination for the study. The study coordinaters anticipated that around 50 participants would be recruited in total by the time of study closure. This was substantially less than the planned (minimum) sample size of 600 participants. This sample size would result in approximately 25 participants allocated to the investigational arm. Therefore, the originally planned analyses may have little value and the models only weakly informed by the available data.

Therefore, this statistical analysis plan was revised to focus the analyses on descriptive summaries of the data and exploratory modelling which might be combined with external information (e.g. historical controls). The originally planned statistical modelling has been retained, but the expectation is that data will be insufficient to reliably undertake these anaylses.

\begin{table}[H]
    \begin{center}
        \begin{tabular}{lllp{6cm}}
            \hline
            Version & Date       & Author           & Description \\
            \hline
            0.1     & 2022-02-08 & James Totterdell & First draft \\
            0.2     & 2022-08-17 & James Totterdell & Revise SAP to focus on descriptive component given early termination of the trial. These revisions were made after termination of the trial, but prior to any data becoming available. \\
            \hline
        \end{tabular}
    \end{center}
\end{table}

\section*{Abbreviations}

\begin{table}[H]
    \begin{center}
        \begin{tabular}{ll}
            \hline 
            Abbreviation & Term \\
            \hline 
            CHS & Clinical Health Score \\
            ICU & Intensive Care Unit \\
            SAP & Statistical Analysis Plan \\
            \hline
        \end{tabular}
    \end{center}
\end{table}

\clearpage

\section{Trial Objectives}

\subsection{Rationale}
The major cause of mortality from COVID-19 has been life-threatening pneumonia and acute respiratory distress syndrome (ARDS). Elevated levels of a number of proinflammatory cytokines, including interleukin-6 (IL-6) and monocyte chemoattractant protein-1 (MCP-1, also known as C-C Motif Chemokine Ligand 2 [CCL2]) have been reported in patients with severe COVID-19, suggesting involvement of a hyper-inflammatory immune response. MCP-1 is a chemokine with a key role in macrophage recruitment and activation and is the natural ligand for Chemokine Receptor Type 2 (CCR2). The utility of CCR2 antagonism in the managing COVID-19 infection has not been tested, with most programs testing CCR2 in the setting of chronic disease such as renal disease.

Alongside the role of chemokines in inducing the hyperinflammatory response, there is good evidence that the renin-angiotensin system (RAS) may also play a role in the pathophysiology of COVID-19. The RAS is responsible for regulating haemodynamic, inflammatory, and fibrotic processes, and includes two main cross-regulating axes: the vasoconstrictive, pro-inflammatory, pro-fibrotic ACE-Ang II-AT1R axis, and the vasodilatory, anti-inflammatory, anti-fibrotic ACE2-Ang1-7-MasR and ACE2-Ang1-9-AT2R axis. SARS-CoV-2 responsible for COVID-19 appears to bind and downregulate angiotensin converting enzyme type 2 (ACE2). Animal studies of the related SARS- CoV-1 suggest that this downregulation of ACE2 is sufficient for causing lung pathology and is reversed by treatment with an Angiotensin Receptor Blocker (ARB). Several randomised controlled trials are underway to assess the effectiveness of ARBs in limiting the severity of COVID-19.

CLARITY 2.0 is an investigator-initiated trial that will evaluate the safety and efficacy of dual treatment with repagermanium, a CCR2 antagonist and candesartan, an ARB, in patients hospitalised with COVID-19 disease.


The full rationale are outlined in the protocol. 


\subsection{Primary Objective and Outcome}
The primary objective is to evaluate the safety and efficacy of dual treatment with repagermanium and candesartan in patients hospitalised with COVID-19 disease, assessed by their \textbf{Clinical Health Score at day 14.}

This Clinical Health Score (CHS) is determined within an 8-point ordinal scale of health status which is a modified version of the 9-point score developed by the WHO for Coronavirus Disease 2019 (COVID-19) trials. A single score will be reported with higher values corresponding to worse symptoms. The ordinal scale is an assessment of the clinical status of the participant at the first assessment for the day, measured at day 14 after the date of randomisation.

The 8-point ordinal scale used for Clinical Health Score is:
\begin{enumerate}[nolistsep]
    \item Not hospitalised, no limitations on activities.
    \item Not hospitalised, limitation on activities.
    \item Hospitalised, not requiring supplemental oxygen.
    \item Hospitalised, requiring supplemental oxygen by mask or nasal prongs.
    \item  Hospitalised, on non-invasive ventilation or high-flow oxygen devices.
    \item Hospitalised, requiring intubation and mechanical ventilation.
    \item Hospitalised, on invasive mechanical ventilation and additional organ support (ECMO).
    \item Death.
\end{enumerate}

\subsection{Secondary Objectives and Outcomes}
The secondary objectives are to evaluate the safety and efficacy of dual treatment with repagermanium and candesartan in patients hospitalised with COVID-19 disease, assessed by:
\begin{enumerate}
    \item Clinical Health Score at day 28.
    \item ICU admission (incidence in days 0-28).
    \item Death (incidence in days 0-28).
    \item Time to death, assessed from hospital admission to death.
    \item Acute Kidney Injury (incidence in days 0-28).
    \item Respiratory Failure (incidence in days 0-28).
    \item Length of hospital admission (days of inpatient stay from admission to discharge or death).
    \item Length of ICU Admission (days in ICU from admission to transfer to ward or death).
    \item Requirement of ventilatory support (count of days with ventilation in days 0-28).
    \item Requirement of dialysis (count of days with dialysis in days 0-28).
    \item Clinical Health Score at day 60.
    \item Clinical Health Score at day 90.
    \item Clinical Health Score at day 180.
\end{enumerate}

\subsection{Safety Objectives and Outcomes}
The specific safety objectives are to evaluate the safety of dual treatment with repagermanium and candesartan in patients hospitalised with COVID-19 disease, assessed by incidence of pre-specified clinical events:

\begin{enumerate}[resume]
    \item Hypotension, requiring an urgent or non-urgent intervention (e.g., reduction in dose or cessation of anti-hypertensive, vasopressors, intravenous fluids). Incidence in days 0-28.
    \item  Hyperkalaemia (defined as a K>5.5-6.0 mmol/L or requiring an intervention including hospitalisation; K>6.0 mmol/L). Incidence in days 0-28.
    \item Deranged Liver Function Tests (defined as Aspartate Aminotransferase (AST) and Alanine Aminotransferase (ALT) >Upper Limit of Normal (ULN) or >1.5 times baseline). Incidence in days 0- 28.
    \item Total SAEs
\end{enumerate}

\clearpage

\section{Study Design}
CLARITY 2.0 is a prospective, Multi-Centre, Multi-Arm Multi-Stage Randomised, Double Blind, Placebo Controlled Trial, utilising adaptive sample size re-estimation. Stage 1 will include 80 patients for a Phase II safety analysis to be conducted in India. Stage 2 will include 520 participants for review of preliminary efficacy data. Expansion to Stage 3, a full Phase III study, will be subject to review of accumulated data in Stage 1 and Stage 2.

Protocol Stage 1 will be conducted in India only. The rest of the world will initiate the study in Stage 2 regardless of the recruitment status of the Stage 1. Recruitment in India will not continue to Stage 2 until completion of the Stage 1 safety analysis and review and approval of the Indian Central Drugs Standard Control Organization Subject Expert Committee on COVID-19 related proposals.


\subsection{Target Population}
Participants must meet all the inclusion criteria, and none of the exclusion criteria, to be eligible for this trial. No exceptions will be made to these eligibility requirements at the time of randomisation. All enquiries about eligibility should be addressed by contacting the sponsor prior to randomisation.

Adults with laboratory-confirmed diagnosis of SARS-CoV-2 infection intended for hospital admission for management of COVID-19.

\subsubsection{Eligibility}
\textbf{Inclusion criteria:}
\begin{itemize}
    \item Adults aged between 18 and 65 years.
    \item Laboratory-confirmed diagnosis of SARS-CoV-2 infection within 10 days prior to randomisation. (Confirmation must be through Reverse Transcription Polymerase Chain Reaction [RT-PCR] method)
    \item Intended for hospital admission for management of COVID-19.
    \item  Patients with moderate (respiratory rate of $\geq$ 24/minute or SPO2: 90\% to $\leq$ 93\% on room air) or severe (respiratory rate of $\geq$ 30/minute or SPO2: <90\% on room air) COVID-19.
    \item  Systolic Blood Pressure (SBP) $\geq$ 120 mmHg OR SBP $\geq$ 115 mmHg and currently treated with a non-RAASi BP lowering agent that can be ceased.
    \item Willing and able to comply with all study requirements, including treatment, timing and/or nature of required assessments.
    \item Documented informed consent
\end{itemize}

\textbf{Exclusion criteria:}
\begin{itemize}
    \item Currently treated with an ACEi, ARB or aldosterone antagonist, aliskiren, or ARNi
    \item Intolerance of ARBs
    \item Serum potassium >5.5 mmol/L
    \item An estimated Glomerular Filtration Rate (eGFR) <30ml/min/1.732m
    \item Known biliary obstruction, known severe hepatic impairment (a Child-Pugh-Turcotte score 10-15)
    \item Pregnancy, lactation, or inadequate contraception.
          \begin{itemize}
              \item Female participants must be either post-menopausal, infertile or use a reliable means of contraception for during the treatment period and for at least 60 days after the last dose of investigational product or refrain. Where they are of childbearing potential, female participants must also have a negative pregnancy test result within 7 days prior to randomisation.
              \item Male participants must have been surgically sterilised or use a (double if required) barrier method of contraception during the treatment period and for at least 60 days after the last dose of investigational product or refrain from donating sperm during this period.
          \end{itemize}
    \item Participation in a study of a novel investigational product within 28 days prior to randomisation.
    \item Plans to participate in another study of a novel investigational product during this study.
\end{itemize}

\subsection{Interventions}
\textbf{Investigational arm (C+R):} Titratable candesartan with commencing dose 4mg tablets twice daily (daily dose 8 mg) + fixed dose repagermanium one x 120mg immediate release capsule twice daily (total daily dose 240mg).

\textbf{Control arm \#1 (C+P):} Titratable candesartan with commencing dose 4mg tablets twice daily (daily dose 8 mg) + matched placebo repagermanium one capsule twice daily.

\textbf{Control arm \#2 (P+P):}  Titratable matched placebo candesartan
one tablet twice daily + matched placebo repagermanium one capsule twice daily.

Treatment for will continue for 28 days.

\emph{Due to early termination of the trial, no participants will have been assigned to the P+P arm.}

\subsection{Randomisation}
For stage 1, treatment allocation will be 1:1 randomisation via block randomisation stratified by centre between the investigational arm (C+R) and control arm 1 (C+P).

Following stage 1, in stage 2 the three arms will be randomised 1:1:1 using block randomisation stratified by centre.

\subsection{Blinding}
The packaging and labelling of interventions will be designed to maintain blinding to the study team and to participants (double-blind).

\subsection{Sample Size}

\emph{Prior to completion of the SAP, the trial was terminated early at the discretion of the funder during recruitment for stage 1 of the trial (planned data collection for safety review). The expected final sample size post-termination was 50 participants. Therefore, the originally planned sample size described below does not apply. It's expected that these 50 participants will be approximately distributed as 25 to each of C+R and C+P given the 1:1 randomisation.}

Stage 1 of the trial will recruit 80 participants from India for a safety analysis.

Stage 2 of the trial will recruit an additional 520 participants. Information from other relevant trials will be utilised to inform the decision for dropping one of the two control arms.

Progression to stage 3 will be subject to review of the combined stage 1 and stage 2 data.

\clearpage

\section{Descriptive Analyses}

\subsection{Participant Throughput}
The flow of participants through the trial will be summarised for each arm using a CONSORT diagram. The flow diagrams will describe the numbers of participants randomly allocated, who received allocation, withdrew consent, and included in the ITT analysis population.

\subsection{Baseline Characteristics}
The following characteristics will be described separately for patients randomised to each arm:
\begin{itemize}
    \item age at randomisation
    \item sex
    \item ethnicity
    \item weight
    \item height
    \item smoking history
    \item comorbidities
    \item concomitant medications
    \item recent blood test results
\end{itemize}
In general, discrete data will be summarised by counts and proportions of participants within each level of the outcome. Continuous variables will be summarised by median, lower quartile, upper quartile, minimum and maximum, or where appropriate, by mean and standard deviation. Counts and proportions of missing baseline values will be reported for each baseline variable.

\subsection{Completeness of Follow-up}
The number and percentage of participants with follow-up information at day 14 and at day 28 post randomisation will be reported. Patterns of missingness will be summarised for primary and secondary outcomes.

\subsection{Treatment Adherence}
The number and proportion of patients who did not receive the treatment they were allocated to will be reported (if any). Concomitant medications received during the treatment period will be reported. Details on the number of days/doses of treatment received will be reported for each arm and compared with the treatment protocol. The number and proportion of participants who discontinued treatment, the timing of discontinuation, and the reasons for discontinuation as specified in the protocol will be reported.

\subsection{Outcomes}

The following planned descriptive summaries outline the minimum required for reporting for the study outcomes. Additional summaries (figures, tables, other summary statistics) may be provided as appropriate or as requested.

For all outcomes, summaries will be grouped by the assigned study arm. For all outcomes, the count and proportion of participants with missing values should be reported.

\subsubsection{Primary Outcome}
Clinical health score is an ordinal outcome with 8 levels. For clinical health score at day 14, the count and proportion of participants in each level of the ordinal scale will be reported by treatment arm along with the count and proportion with missing outcome.

\subsubsection{Secondary Outcomes}

The following should at least be reported for each secondary outcome:

\begin{enumerate}
    \item \textbf{Clinical Health Score at day 28 (ordinal):} \\
        As per the primary outcome.
    \item \textbf{ICU admission days 0 to 28 (binary):} \\
        ICU admission is a binary outcome (assuming that a patient can only have one ICU admission within 28 days). The count and proportion of participants who had an ICU admission. Additionally, a cross tabulation of ICU admission with death: neither death nor ICU, ICU without death, death without ICU admission, ICU followed by death.
    \item \textbf{Death days 0 to 28 (binary):} \\
        The count and proportion of participants who died within 28 days.
    \item \textbf{Time to death (days), assessed from hospital admission (discrete):} \\
        Time to death will be censored at either the end of follow-up (180 days post-randomisation) or in the case of loss-to-follow-up, at the last study day on which the participant was known to be alive. Kaplan-Meier curves (or a discrete-time equivalent) for time to death will be presented.
    \item \textbf{Acute Kidney Injury days 0 to 28 (binary):} \\
        The count and proportion of participants who experienced AKI within 28 days.
    \item \textbf{Respiratory Failure days 0 to 28 (binary):} \\
        The count and proportion of participants who experienced RF within 28 days.
    \item \textbf{Length of hospital admission, days of inpatient stay from admission to discharge or death (discrete):}
        The distribution of days in hospital and the count/proportion of participants with each value. The median, Q1, and Q3 days in hospital.
    \item \textbf{Length of ICU admission, days in ICU from admission to transfer to ward or death, (discrete):} \\
        The distribution of days in ICU and the count/proportion of participants with each value. The median, Q1, and Q3 days in ICU amongst all participants and amongst those who spent any time in ICU.
    \item \textbf{Requirement of ventilatory support, days with ventilation on days 0 to 28 (discrete):} \\
        The count and proportion of participants who required any ventilatory support. The distribution of number of days requiring ventilatory support and relevant summaries of this distribution (median, Q1, Q3 etc).
    \item \textbf{Requirement of dialysis, days with dialysis in days 0 to 28 (discrete):} \\
        The count and proportion of participants who required any dialysis. The distribution of number of days requiering dialysis and relevant summaries of this distribution (median, Q1, Q3 etc).
    \item \textbf{Clinical Health Score at day 60:} \\
        As per the primary outcome.
    \item \textbf{Clinical Health Score at day 90:} \\
        As per the primary outcome.
    \item \textbf{Clinical Health Score at day 180:} \\
        As per the primary outcome.
\end{enumerate}

\subsubsection{Safety Outcomes}

The following should at least be reported for each secondary outcome:

\begin{enumerate}[resume]
    \item \textbf{Hypotension, requiring an urgent or non-urgent intervention (e.g., reduction in dose or cessation of anti-hypertensive, vasopressors, intravenous fluids) days 0-28 (binary):} \\
        The count and proportion of participants who met the definition within 28 days.
    \item  \textbf{Hyperkalaemia (defined as a K>5.5-6.0 mmol/L or requiring an intervention including hospitalisation; K>6.0 mmol/L) days 0-28 (binary):} \\ 
        The count and proportion of participants who met the definition within 28 days.
    \item \textbf{Deranged Liver Function Tests (defined as Aspartate Aminotransferase (AST) and Alanine Aminotransferase (ALT) >Upper Limit of Normal (ULN) or >1.5 times baseline) days 0- 28 (binary):} \\
        The count and proportion of participants who met the definition within 28 days.
    \item \textbf{Total SAEs}:
        The count and proportion of participants who experienced any SAEs. The total number and rate of any SAEs.
\end{enumerate}

\subsubsection{Additional Summaries}

If available, the daily distribution (between day 0 to 28) of the clinical health score outcome will be summarised (counts and proportions) for each treatment group and may be presented as a stacked bar chart with day on the x-axis and cumulative proportion on the y-axis, coloured by the ordinal levels.

\clearpage

\section{Comparative Analyses}

\textit{The following comparative statistical analyses were originally planned for the study outcomes. However, due to the early trial termination and small expected sample size the originally planned analyses may have little value or be only weakly informed from the data due to small numbers of events (e.g. time to death, length of ICU stay, etc.). For this reason, the planned analyses introduced in the following may not be undertaken or may be greatly simplified at the discretion of the analyst.}

\textit{Alternatively, the analyses described (or similar) may be used as an exploratory venture where it's considered to be of value, or used as a guide for models considered in investigating methods which incorporate historical controls (from CLARITY 1.0) or other external information into the analyses in a more considered way. In the absence of any pre-specified use of external controls, any analyses incorporating them should be identified as exploratory.}

\textit{The planned analyses were for 3 study arms. Given no participants were allocated to the double placebo arm, one of the intervention covariates ($\beta$) should be dropped from each model.}

Pairwise comparisons will be made between the investigational arm and each control arm. For all outcomes, the primary analyses will be performed with treatment groups as assigned by randomisation (de facto estimand) irrespective of treatment withdrawal or failure to comply with the protocol.

For all models, any baseline covariates included for adjustment will be mean-centred where appropriate, and the treatment design (denoted by $x$) will use treatment contrasts such that intercepts represent the outcome distribution amongst patients assuming equal weighting across all treatment groups.

The posterior distribution of each contrast of interest will be reported along with posterior summaries: median, 95\% highest density credible interval, and posterior probability of events of interest (e.g. that $\beta < 0$). Posterior summaries of other model parameters will be reported.

\subsection{Intercurrent Events}

For most analyses, the only intercurrent event (ICE) of concern is death. In most cases, this ICE is handled by the use of composite outcomes where death is assigned the worst possible score of the outcome on the ordinal scale.

For other intercurrent events (e.g. treatment switching/withdrawal) the treatment-policy strategy will be used with treatment groups analysed as assigned irrespective of withdrawal from or discontinuation of the assigned treatment.

\subsection{Primary Outcome}\label{sec:primary-outcome-analysis}
The following Bayesian cumulative logistic regression model will be updated using the available data.
$$
    \begin{aligned}
        y_{i} | \pi;x_i                            & \sim \text{Categorical}(\pi(x_i))                              \\
        \mathbb P(y_i \leq k | \alpha, \beta; x_i) & = \text{logit}^{-1}(\eta_{ik})                                 \\
        \eta_{i}                                   & = x_i^{\mathsf{T}}\beta \\
        \pi_k(\eta)                                & = \begin{cases}
            \text{logit}^{-1}(\alpha_1 - \eta)                                            & k=1             \\
            \text{logit}^{-1}(\alpha_{k} - \eta) - \text{logit}^{-1}(\alpha_{k-1} - \eta) & k\in\{2,...,7\} \\
            1 - \text{logit}^{-1}(\alpha_{7} - \eta)                                      & k=8
        \end{cases},
    \end{aligned}
$$
with prior distributions
$$
    \begin{aligned}
        \pi(0)          & \sim \text{Dirichlet}(\kappa)                            \\
        \beta_1,\beta_2 & \sim \text{Normal}(0, 1)                                 \\
        \kappa          & = 4\cdot(0.80, 0.11, 0.02, 0.02, 0.01, 0.01, 0.01, 0.02)
    \end{aligned}
$$
where the hyper-parameters $\kappa$ are weakly informed by the observed distribution at day 14 in CLARITY 1. The weight of 4 given to this prior was chosen ad-hoc for model stability.

As a prior sensitivity check, the model should also be fit assuming a weakly informative prior on the outcome levels, i.e.
$$
\kappa = 2 \times (1/8, 1/8, 1/8, 1/8, 1/8, 1/8, 1/8, 1/8)
$$

Due to the planned use of orthornormal design coding, the prior specified on $\pi(0)$ is the prior distribution on the outcome levels for an average patient with equal weighting given to each study arm, rather than the distribution of outcome levels for patients in the double placebo control group (\textit{note that due to early termination, no participants will have been assigned to the double placebo control group, so $\beta_2$ should be dropped from the model and it may be that treatment coding are more appropriate for the design with control arm \#1 as the reference group}).

\subsection{Secondary Outcomes}

\subsubsection{Clinical health score at day 28}
The analysis will be analagous to that for clinical health score at day 14 (\ref{sec:primary-outcome-analysis}), but with a different prior. The prior is weakly informed by CLARITY 1 outcome data at day 28. For this outcome we specify
$$
    \kappa = 4\times(0.91, 0.02, 0.01, 0.01, 0.01, 0.01, 0.01, 0.02)
$$
for the hyper-parameter on the prior for $\pi(0)$.

\subsubsection{ICU admission (incidence days 0 - 28)}\label{sec:icu-analysis}
The following Bayesian logistic regression model will be updated using the available data.
$$
    \begin{aligned}
        y_i|\pi_i & \sim \text{Bernoulli}(\pi_i)   \\
        \pi_i     & = \text{logit}^{-1}(\eta_i)    \\
        \eta_i    & = \alpha + x_i^\mathsf{T}\beta
    \end{aligned}
$$
with priors
$$
    \begin{aligned}
        \alpha & \sim \text{Logistic}(\text{logit}(0.2), 1)     \\
        \beta  & \overset{\text{iid}}{\sim} \text{Normal}(0, 1)
    \end{aligned}
$$
where the prior for $\alpha$ is weakly informed by CLARITY 1 results.

The model above ignores deaths (i.e. participants who died without requiring ICU admission are counted as no ICU admission). If it is necessary to account for deaths as part of the outcome, then the outcome may be modified to be either:
\begin{itemize}
    \item ICU admission or death vs neither (binary)
    \item death (2), ICU admission (1), neither (0) (ordinal)
\end{itemize}
with similar priors weakly informed from CLARITY 1.0.

\subsubsection{Death (incidence days 0 - 28)}
This outcome will be analysed analgously to the ICU incidence outcome (\ref{sec:icu-analysis}). However, the following priors will instead be specified
$$
    \alpha \sim \text{Logistic}(\text{logit}(0.1), 1),
$$
where the prior is weakly informed by CLARITY 1 results.

\subsubsection{Time to death (from randomisation)}
Time to death will be analysed assuming a proportional hazards model. The baseline hazard will be modelled flexibly using M-splines.

Let $D_i = (y_i, \nu_i)$ denote the data for subject $i$ where $y_i$ is the event time and $\nu_i$ and indicator for censoring or observation of the event. The model for the hazard is
$$
    \begin{aligned}
        \lambda(t|\beta;x) & = \lambda_0(t)\exp(x^\mathsf{T}\beta)                           \\
        \lambda_0(t)       & = \sum_{l=1}^L M_l(t;\tau)\gamma_l,\quad \sum_{l=1}^L\gamma_l=1 \\
        \eta_i             & = x_i^\mathsf{T}\beta                                           \\
        \lambda_i(y_i)     & = \lambda_0(y_i)\exp(\eta_i)                                    \\
    \end{aligned}
$$
where $M_l(t;\tau)$ is the $l$th basis term for a M-spline of degree 3 with knots at locations $\tau=\{\tau_1,...,\tau_{L+4}\}$ evaluated at $t$. The default not location will be at the quantiles of observed event times. If this is determined to be inappropriate (too few events, insufficient flexible etc.) then an alternative may be specified (e.g. parametric baseline hazard).

The priors are
$$
\begin{aligned}
    \gamma & \sim \text{Dirichlet}(\kappa) \\
    \beta  & \overset{\text{iid}}{\sim} \text{Normal}(0, 1)
\end{aligned}
$$
where $\kappa = 2\times\{1/L\}_{l=1}^L$.

The contrasts of interest will be the hazard ratio of the intervention group relative to each of the control groups.

If there is clear departure from proportional hazards assumption, then the treatment effect parameters $\beta$ may be allowed to vary with time rather than a time-invariant effect. In this case, restricted mean survival time will be the comparison of interest.

\subsubsection{Acute kidney injury (incidence in days 0 - 28)}
This outcome will be analysed analgously to the ICU incidence outcome (\ref{sec:icu-analysis}). However, the following priors will instead be specified
$$
    \alpha \sim \text{Logistic}(\text{logit}(0.2), 1),
$$
where the prior is weakly informed by CLARITY 1 results.

\subsubsection{Respiratory failure (incidence in days 0 - 28)}
This outcome will be analysed analgously to the ICU incidence outcome (\ref{sec:icu-analysis}). However, the following priors will instead be specified
$$
    \alpha \sim \text{Logistic}(\text{logit}(0.2), 1),
$$
where the prior is weakly informed by CLARITY 1 results.

\subsubsection{Length of hospital admission (days of inpatient stay from admission to discharge or death)}

This outcome is equivalently stated as the time to hospital discharge (either due to discharge home or death). Death should be treated as a competing event to being discharged alive. For simplicity (given that most participants will have died or been discharged within 28 days based in previous studies), the time to event may be treated as censored at day 28 and the outcome analysed by a time-to-event model for the competing events of discharge alive or death to day 28 assuming constant hazard associated with the intervention across all days from 1 to 28.

For a discrete time-to-event model, denote by $d_i\in\{1,...,28\}$ denote the day of first event (or censoring at day 28) and $m\in\{0,1,2\}$ indicating the event type (still admitted to hospital, discharged alive, died) and $y_{id} = (y_{id0}, y_{id1}, y_{id2})$ with $y_{itm}\in\{0,1\}$ indicates the participant $i$ status at day $d\leq d_i$. Then
$$
\begin{aligned}
\lambda_m(d, \eta) &= \frac{\exp(\alpha_{dm} + \eta_m)}{1 + \sum_{l=1}^M \exp(\alpha_{dl} + \eta_l)},\quad m=1,2 \\
y_{id} &\sim \text{Multinomial}\left(1, \left(1 - \sum_{l=1}^2 \lambda_i(d,\eta_i), \lambda_1(d, \eta_i), \lambda_2(d, \eta_i)\right)\right),\quad d\leq d_i
\end{aligned}
$$
where $\eta=(\eta_1,\eta_2)$ contains the event-specific linear predictors with event specific parameters.

\subsubsection{Length of ICU admission (days in ICU from admission to transfer to ward or death)}

\textit{No analysis plan for this outcome had been written prior to early termination of the trial. Not all participants will have an ICU admission, this outcome will only be relevant to participants admitted to ICU. Given the small expected sample size (50) and the expectation that there will be a small number of ICU admissions, no analysis is planned beyond the descriptive summaries previously described.}

\subsubsection{Requirement of ventilatory support (number of days with ventilation in days 0 - 28)}\label{sec:vent-support}

A cumulative logistic proportional odds model will be assumed. If $y_i\in\{0,1,...,28,29\}$ denotes the number of days with ventilatory support (where 0 means 0 days, and 29 corresponds to death) for participant $i$, and $x_i$ the intervention design covariates, then
$$
    \begin{aligned}
        y_{i} | \pi;x_i                            & \sim \text{Categorical}(\pi(x_i))           \\
        \mathbb P(y_i \leq k | \alpha, \beta; x_i) & = \text{logit}^{-1}(\eta_{ik})              \\
        \eta_{ik}                                  & = \alpha_k + x_i^{\mathsf{T}}\beta + \cdots \\
        \pi_k(x)                                   & = \begin{cases}
            1 - \text{logit}^{-1}(\alpha_1 + x^\mathsf{T}\beta)                                                     & k=0               \\
            \text{logit}^{-1}(\alpha_{k-1} + x^\mathsf{T}\beta) - \text{logit}^{-1}(\alpha_{k} + x^\mathsf{T}\beta) & k\in\{2,...,K-1\} \\
            \text{logit}^{-1}(\alpha_{K-1} + x^\mathsf{T}\beta)                                                     & k=29
        \end{cases},
    \end{aligned}
$$
with prior distributions
$$
    \begin{aligned}
        \pi(0)          & \sim \text{Dirichlet}(\kappa)  \\
        \beta_1,\beta_2 & \sim \text{Normal}(0, 1).
    \end{aligned}
$$
where $\kappa = 2\times \{1/30\}_{k=0}^{29}$.

The primary contrasts will be the relative log-odds of having a worse outcome (higher outcome level) in the investigational arm compared to each of the contorl arms.

If the outcome is sparse across the levels, then outcome levels may be aggregated (e.g. 0 days of ventilation vs any days ventilation in the most extreme case). The prior should be adjusted appropriately. 

Non-proportionality, particularly with respect to death, to be investiated.

\subsubsection{Requirement of dialysis (number of days with dialysis in days 0 - 28)}

This outcome will be analysed analgously to \ref{sec:vent-support}.

\subsubsection{Clinical health score at day 60}
The analysis will be analagous to that for clinical health score at day 14 (\ref{sec:primary-outcome-analysis}), but with a different prior. The prior is weakly informed by CLARITY 1 outcome data. For this outcome we specify
$$
    \kappa = 2\times(0.91, 0.02, 0.01, 0.01, 0.01, 0.01, 0.01, 0.02)
$$
for the hyper-parameter on the prior for $\pi(0)$.

\subsubsection{Clinical health score at day 90}
The analysis will be analagous to that for clinical health score at day 14 (\ref{sec:primary-outcome-analysis}), but with a different prior. The prior is weakly informed by CLARITY 1 outcome data. For this outcome we specify
$$
    \kappa  = 2\times(0.91, 0.02, 0.01, 0.01, 0.01, 0.01, 0.01, 0.02)
$$
for the hyper-parameter on the prior for $\pi(0)$.

\subsubsection{Clinical health score at day 180}
The analysis will be analagous to that for clinical health score at day 14 (\ref{sec:primary-outcome-analysis}), but with a different prior. The prior is weakly informed by CLARITY 1 outcome data. For this outcome we specify
$$
    \kappa = 2\times(0.91, 0.02, 0.01, 0.01, 0.01, 0.01, 0.01, 0.02)
$$
for the hyper-parameter on the prior for $\pi(0)$.

\subsection{Safety Outcomes}

For safety outcomes no formal modelling is planned beyond the descriptive statistics previously described. If required, logistic regression or count regression models with weakly informative priors may be used as relevant to the safety outcome.

\subsection{Baseline Adjustments}

\textit{Due to the small expected sample size, it's suggested that no baseline covariate adjustments are made. The following reports on the originally planned adjustments.}

As randomisation is stratified by centre, all models will include centre as a random effect. The effect of centre will be assumed additive in the linear predictor and will be modelled by
$$
    \begin{aligned}
        \gamma_j|\tau & \sim \text{Normal}(0, \tau^2), \quad j =1,...,J \\
        \tau          & \sim \text{Cauchy}(0, 2.5).
    \end{aligned}
$$

\subsection{Subgroup Analyses}
No pre-specified subgroup analyses are planned.

\subsection{Missing Data}
Patterns of missing baseline and outcome data will be reported by treatment group.

Where other sources of information are available, missing baseline variables may be imputed deterministically. Missing outcomes may be treated as censored if appropriate (e.g. if a participant was known to be alive at day 14, but their exact status was unknown then their outcome is interval censored in $\{1,...,7\}$) or if they were known to be out of hospital but with unknown not-hospitalised status then their outcome is interval censored in $\{1,2\}$).

Baseline predictors of missingness will be investigated. In the absence of strong predictors of missingness, the default analysis will be based on available cases adjusting for the pre-specified covariates. Alternatively, missing outcomes may be multiply imputed using an expanded set of baseline covariates and the combined results used for reporting.

\subsection{Software}
The statistical software R will be used for data processing. Bayesian model posteriors will be approximated using HMC as implemented in Stan. Program and package versions used for the analyses will be reported.

\clearpage

\section{Trial Monitoring and Reporting}

\textit{Due to early trial termination there will be no trial monitoring. The following is retained for reference to the original plan.}

Analyses are planned to occur at the following stages of recruitment:

\begin{itemize}
    \item Exploratory safety analysis after 80 particiants randomsied 1:1 to two arms.
    \item Effectiveness analysis after an additional 520 participants randomised 1:1:1 to all 3 arms.
    \item Effectiveness analysis after every additional 600 participants up to the maximum sample size of 2,100 participants.
\end{itemize}

At each analysis, the investigational arm will be compared with both control arms with respect to the primary outcome.

\subsection{Decision Rules}

Decisions pertaining to continuation of the trial will be guided by the calculation of predictive probabilities of satisfying the pre-specified criteria of effectiveness with respect to the primary outcome.

The quantities of interest are the predictive probability of trial success for each of the contrasts, defined as
\begin{align}
    \text{PPoS}_{\text{P+P}}(\texttt{data}_n, m) & = \mathbb{E}[\mathbb{P}(\eta_3 - \eta_2 > \delta | \texttt{data}_{n+m}) > \epsilon_{\text{eff}} | \texttt{data}_n] \\
    \text{PPoS}_{\text{C+P}}(\texttt{data}_n, m) & = \mathbb{E}[\mathbb{P}(\eta_3 - \eta_1 > \delta | \texttt{data}_{n+m}) > \epsilon_{\text{eff}} | \texttt{data}_n]
\end{align}
where (1) relates to the comparison of the investigational arm with control arm 1 and (2) the comparison of the investigational arm with control arm 2.

Given the aim of the trial is to establish effectiveness relative to both of the control arms, if either of the quantities falls below a threshold, $\epsilon_{\text{fut}}$, at an interim analysis then a decision of trial futility is recommended.

Calculation of these predictive quantities requires assumptions about the future distribution of covariates included in the primary model. The assumption will be that future participants have similar covariates to past participants.

\clearpage

\printbibliography[heading=bibintoc]

\end{document}
